\chapter{Introduction to Deep Learning}
\label{ch:01-intro-deep-learning}

\lecturemeta{Sarath Chandar}{Polytechnique Montr\'{e}al and Mila}

The lecture is organised around a single asymmetry: we have known since the early 1990s
what a neural network \emph{can} represent, and we still do not know what we can reliably
\emph{train}. Chandar spends the first half establishing the representational result, the
universal approximation theorem, in order to set it aside, and the second half on the
obstructions that survive it: gradients that vanish through depth, initialisations that
decide whether a network trains at all, and finally the independence assumption buried in
every objective on the slides. The closing frame is the argument of the whole summer school
in one line: current practice optimises the i.i.d.\ case, while ``true learning'' is
sequential, online and continual.

\section*{Key ideas}

\paragraph{The gradient that cancels.} The derivation that opens the lecture is a single
sigmoid unit, $y^{(n)} = \sigmoid(a^{(n)})$ with $a^{(n)} = \wt\T\xin^{(n)}$, trained on
the cross-entropy objective
\[
  \min_{\wt} \; -\sum_{n=1}^{N} \Bigl[ t^{(n)} \ln y^{(n)}
    + \bigl(1 - t^{(n)}\bigr) \ln\bigl(1 - y^{(n)}\bigr) \Bigr],
  \qquad
  \frac{\partial \mathrm{obj}}{\partial \wt}
   = \sum_{n=1}^{N} \frac{\partial \mathrm{obj}}{\partial y^{(n)}}
     \frac{\partial y^{(n)}}{\partial a^{(n)}}
     \frac{\partial a^{(n)}}{\partial \wt}.
\]
\noindent Worked factor by factor, the objective contributes
$\partial \mathrm{obj} / \partial y^{(n)} = (y^{(n)} - t^{(n)}) / (y^{(n)}(1 - y^{(n)}))$,
the unit contributes $\partial y^{(n)} / \partial a^{(n)} = y^{(n)}(1 - y^{(n)})$, and the
pre-activation contributes $\partial a^{(n)} / \partial \wt = \xin^{(n)}$. The awkward
denominator is exactly the sigmoid derivative, so the two cancel, as the board shows by
striking them out, and each term of the sum collapses to
\[
  \bigl(y^{(n)} - t^{(n)}\bigr)\xin^{(n)}.
\]
\noindent Residual times input. The cancellation is the point: it is why this pairing of
output unit and loss is the one everybody uses, and it is the atom that backpropagation
chains.

\paragraph{Backpropagation has a messy provenance.} The slides are unusually careful here, and they
begin further back than most treatments: the perceptron in 1957~\citep{rosenblatt1958}, logistic
regression in 1958~\citep{cox1958}, and the XOR counterexample in 1969~\citep{minsky1969}.
Backpropagation was then popularised by Rumelhart, Hinton and Williams in
1986~\citep{rumelhart1986}, but Werbos had it in 1974~\citep{werbos1974}, and the slide
states flatly that it ``has been reinvented several times before 1986'', pointing at
Schmidhuber's history page~\citep{schmidhuber2014backprop}. What actually changed the field
was not the rule but the tooling, and Chandar gives the timeline: Torch (2002), Theano
(2007 to 2008), Caffe (2013), TensorFlow (2015), PyTorch (2016) and JAX (2018). Automatic
differentiation, not the chain rule, is the enabling technology.

\paragraph{Expressivity is settled.} The universal approximation theorem is quoted as
Hornik, 1991~\citep{hornik1991}: a single hidden layer network with a linear output unit can
approximate any continuous function arbitrarily well, \emph{given enough hidden units}. On
the slide the final clause is highlighted and braced, and underneath, in red, sits the
phrase that reframes everything after it: \textbf{expressivity versus learnability}. One
hidden layer suffices in principle, which is precisely why the rest of the lecture is about
optimisation rather than architecture.

\paragraph{Learnability is not.} Depth turns the chain rule into a product of Jacobians,
\[
  \frac{\partial E}{\partial \mat{W}^{(T)}}
   = \frac{\partial E}{\partial \yout}\,
     \frac{\partial \yout}{\partial \vect{h}_T}
     \underbrace{\left[
       \frac{\partial \vect{h}_T}{\partial \vect{h}_{T-1}}
       \frac{\partial \vect{h}_{T-1}}{\partial \vect{h}_{T-2}} \cdots
       \frac{\partial \vect{h}_2}{\partial \vect{h}_1}
     \right]}_{\text{vanishing gradient problem}},
\]
\noindent and the bracket is boxed on the board. The historical answer was to avoid
training the depth all at once: greedy layerwise pretraining~\citep{hinton2006, bengio2007},
in which each layer is fitted as an autoencoder before the stack is finetuned end to end
(Figure~\ref{fig:intro-deep-learning-1}). The modern answer is quieter and lives in the
activation function, where the lecture records
$\mathrm{Swish}(x) = x\,\sigmoid(x)$~\citep{ramachandran2017} and
$\mathrm{SwiGLU}(x) = \mathrm{Swish}(x\mat{W}) \otimes x\mat{V}$~\citep{shazeer2020}.

\begin{figure}[htbp]
\centering
\begin{tikzpicture}[
  x=1cm, y=1cm, font=\small,
  box/.style={draw, rounded corners=1pt, minimum height=8mm, minimum width=7mm,
              inner sep=1pt, fill=white},
  hid/.style={box, fill=asidebg},
  ar/.style={-{Latex[length=1.6mm]}, shorten >=1pt, shorten <=1pt}]

  \node[box] (x1) at (0,0) {$\xin$};
  \node[hid] (h1) at (1.6,0) {$\vect{h}_1$};
  \node[box] (xr) at (3.2,0) {$\est{\xin}$};
  \draw[ar] (x1) -- node[above,font=\scriptsize] {$\mat{W}^{(1)}$} (h1);
  \draw[ar] (h1) -- node[above,font=\scriptsize] {$\mat{W}^{(1)\mathsf{T}}$} (xr);
  \node[font=\scriptsize, align=center] at (1.6,-1.05)
    {train layer 1 as an autoencoder,\\ reconstructing the input};

  \node[box] (x2) at (6.3,0) {$\xin$};
  \node[hid] (h1b) at (7.9,0) {$\vect{h}_1$};
  \node[hid] (h2) at (9.5,0) {$\vect{h}_2$};
  \node[box] (hr) at (11.1,0) {$\est{\vect{h}}_1$};
  \draw[ar] (x2) -- (h1b);
  \draw[ar] (h1b) -- node[above,font=\scriptsize] {$\mat{W}^{(2)}$} (h2);
  \draw[ar] (h2) -- (hr);
  \node[font=\scriptsize, align=center] at (8.7,-1.05)
    {repeat upward: each new layer\\ reconstructs the one below};

  \node[box] (x3) at (2.2,-3.0) {$\xin$};
  \node[hid] (s1) at (3.8,-3.0) {$\vect{h}_1$};
  \node[hid] (s2) at (5.4,-3.0) {$\vect{h}_2$};
  \node (dots) at (6.7,-3.0) {$\cdots$};
  \node[hid] (sL) at (8.0,-3.0) {$\vect{h}_L$};
  \node[box] (y3) at (9.6,-3.0) {$\yout$};
  \draw[ar] (x3) -- (s1); \draw[ar] (s1) -- (s2); \draw[ar] (s2) -- (dots);
  \draw[ar] (dots) -- (sL); \draw[ar] (sL) -- (y3);
  \draw[decorate, decoration={brace, amplitude=4pt, mirror}]
    (1.75,-3.6) -- (10.05,-3.6)
    node[midway, below, yshift=-3pt, font=\scriptsize] {supervised finetuning};
\end{tikzpicture}
\caption{Greedy layerwise pretraining as drawn in the Introduction to Deep Learning
lecture. Each layer is first trained to reconstruct the representation below it, so no
gradient ever has to travel the full depth, and only then is the assembled stack finetuned
against labels. Redrawn from photographs of the board.}
\label{fig:intro-deep-learning-1}
\end{figure}

\paragraph{Initialisation decides trainability.} The lottery ticket
hypothesis~\citep{frankle2019} is quoted verbatim: dense, randomly initialised feed-forward
networks contain subnetworks, ``winning tickets'', that when trained in isolation reach test
accuracy comparable to the original network in a similar number of iterations. Placed here,
immediately after the vanishing-gradient material, it reads as a statement about
optimisation rather than about compression: most of the network is scaffolding for finding a
subnetwork the gradient can actually move.

\paragraph{The i.i.d.\ assumption is the frontier.} The penultimate slide is four lines.
Current focus: i.i.d.\ learning. True learning: non-i.i.d., namely sequential, online and
continual. Nothing follows it except a pointer to Mark Schmidt's optimisation tutorial at
ICML 2026. The lecture ends by naming its own limitation.

\section*{Notation}

Superscript $(n)$ indexes training examples, so $t^{(n)}$ is the target and $y^{(n)}$ the
prediction for example $n$. Layer index is a superscript in parentheses on weight matrices,
$\mat{W}^{(\ell)}$, and a subscript on activations, $\vect{h}_\ell$. This convention
holds for later chapters, with one difference worth flagging now:
Chandar writes the objective as an average over $N$ examples with an explicit minus sign,
whereas Cho (Chapter~\ref{ch:05-statistics-estimation}) works with expectations under a
population distribution. The quantities are the same; the framing differs.

\section*{Why it matters}

This chapter is the hinge for three later ones. The non-i.i.d.\ closing slide is picked up
directly by Lyle on continual learning (Chapter~\ref{ch:06-continual-learning}), who turns
``sequential, online, continual'' into two concrete failure modes, and by Pascanu
(Chapter~\ref{ch:07-to-be-figured-out}), who argues that the i.i.d.\ assumption is what
makes gradient-based credit assignment work at all. The lottery ticket hypothesis reappears
in Alistarh's history of compression (Chapter~\ref{ch:10-efficient-learning}), filed under
pruning rather than optimisation, which is a real difference of emphasis between two
speakers citing one paper.

\begin{aside}
The two readings of lottery tickets pull apart. Chandar places the result among
answers to ``how do we train deep networks'', making it a claim about which initialisations
admit a trainable path. Alistarh places it in the ``Age of Discovery'' of model compression,
making it a claim about how much of a trained network is redundant. The paper supports both,
but they suggest different experiments.
\end{aside}
